\begin{figure*}
    \small
    \centering\includegraphics[width=0.7\linewidth]{figures/maya_simulator_anatomy.pdf}
    \caption{Comparison of modeling approaches: Traditional systems require explicit workload specification and several complex heuristics steps to obtain a kernel level execution graph apt for simulation. On the other hand, \sysname directly captures the computation graph at a lower-level through transparent device emulation.}
    \label{fig:simanatomy}
\end{figure*}


\begin{figure}[htbp]
    \centering
    \includegraphics[width=0.55\linewidth]{figures/maya_hld.pdf}
    \caption{\sysname architecture: (1) Unmodified training code is executed through a virtual runtime that emulates the device drivers given emulation specs, (2) Worker traces capturing device API calls are merged into a unified trace, (3) Kernels in the unified trace are annotated with predicted runtimes using pre-trained estimators, (4) An event-driven simulator processes the annotated traces using cluster specifications to produce performance predictions.}
    \label{fig:hld}
\end{figure}


\vspace{-0.5em}

\section{\sysname : System Design}\label{sec:design}

\begin{figure*}
    \centering
    \includegraphics[width=\linewidth]{figures/trace_flow.pdf}
    \caption{\small \sysname's trace processing pipeline: Starting with unmodified user code, the device emulator captures raw traces containing API calls, kernel launches, and synchronization events across multiple GPU streams. The trace collator merges these per-GPU traces and resolves collective operations, creating a unified job-level trace. The kernel runtime estimator then annotates compute operations with predicted durations. Finally, the event-driven simulator processes this trace to model the complex interactions between compute operations, synchronization events, and communication collectives across streams and devices, producing an accurate timeline of execution.}
    \label{fig:traceflow}
\end{figure*}


The foundation of \sysname is a transparent device emulator that functions as the interface between unmodified training workloads and the modeling pipeline. This component interposes on the accelerator device APIs, virtualizing device interactions while maintaining execution fidelity. The emulator captures a precise trace of device operations --- compute kernels, memory operations, and synchronization events --- completely on the CPU. This yields detailed execution traces while remaining fully transparent to the training application, which executes as if on a real accelerator.

These raw execution traces then flow through a trace collation and analysis pipeline that reconstructs the distributed execution pattern. The collator combines traces from multiple workers, resolving dependencies across both space (between workers) and time (within execution streams) by identifying collective communication operations, which are crucial for modeling distributed training workloads. The performance estimation phase then augments this execution trace with runtime predictions. Since the emulator captures operation metadata but does not execute compute kernels, \sysname employs a combination of machine learning and analytical models to predict operation runtimes.

The final phase uses event-driven simulation to model end-to-end execution. The simulator processes the annotated trace according to a specified hardware configuration, modeling complex execution dependencies within and across workers in a distributed training workload. This captures critical performance characteristics like pipeline bubbles and compute-communication overlap that emerge from the interaction between device operations. The output is a comprehensive simulation report that encompasses metrics such as batch execution time, communication time \& memory usage.

\sysname's architecture enables it to capture the full complexity of modern DLT optimizations while providing high-fidelity performance predictions. By operating on unmodified user code and eliminating the requirement for accelerator hardware during prediction, \sysname offers a unique combination of transparency and efficiency. In the rest of this section, we provide the design details of each component in \sysname (Figure \ref{fig:hld}).

\subsection{Transparent Accelerator Emulation}
We make a key observation on the nature of deep learning training (DLT) workloads: the CPU-side control flow of the application is fundamentally decoupled from the computation executed on accelerator devices. Since there is minimal feedback to the control flow from the results of device operations, it is possible to emulate device behavior without materializing output values. \sysname's emulator exploits this characteristic --- turning compute operations into no-ops while carefully managing device state and dependencies.


To achieve transparency, the emulator \textit{intercepts} calls to device APIs without requiring modifications to the training application. Most device operations, particularly compute kernels, are transformed into no-ops that record metadata about the operation and return immediately. However, \sysname must still precisely track  device state to ensure correct execution. For instance, when applications query device state through APIs like \texttt{cudaMemGetInfo}, \sysname returns carefully constructed responses that mimic device behavior, allowing frameworks like PyTorch to make memory management decisions as they would on real hardware.

The design of this architecture addresses three key challenges in maintaining execution fidelity despite no-op execution: (1) maintaining the semantic meaning of API sequences, (2) tracking both physical and virtual resources, and (3) handling distributed dependencies. These challenges guide \sysname's semantically-aware emulation layer.

\vheading{Context-aware Operation Modeling} Several device operations gain meaning only when considered within the context of a broader sequence of API calls, requiring careful state tracking. In CUDA, for instance, \texttt{cudaStreamWaitEvent()} synchronizes ops across compute streams based on events recorded by \texttt{cudaEventRecord()}. We maintain a map of device state to model these dependencies correctly, even though the underlying operations don't execute.

A similar treatment is required for operations involving opaque libraries like cuBLAS and cuDNN, where configurations are built incrementally. For instance, a cuBLAS matrix multiplication involves a sequence of setup calls (\texttt{cublasSetMatrix()}, \texttt{cublasSetStream()}) before the actual computation (\texttt{cublasGemmEx()}). \sysname tracks these stateful API sequences to construct the complete operation metadata, essential for modeling performance-critical operations like matrix multiplications and convolutions.

\vheading{Resource Tracking} \sysname maintains a dynamic mapping of both physical and virtual resources during emulation. For memory management, \sysname tracks allocations and deallocations, allowing it to simulate real hardware constraints and detect errors such as out-of-memory (OOM) conditions and invalid memory accesses. In unified memory configurations, \sysname tracks tensor locations across host and device spaces and resolves ambiguity in API calls like \texttt{cudaMemcpyAsync} to accurately model workload behavior. 

In addition to physical resources, \sysname creates and manages virtual resources and handles that are returned to the application; examples of this include device handles, CUDA streams, and CUDA events. Any misconfiguration or user error --- such as using an invalid stream or an uninitialized descriptor --- is identified and flagged by \sysname using each handle's state. Through detailed accounting of both physical and virtual resources, \sysname provides a realistic foundation to emulate device behavior and potential failure scenarios; this is a key benefit unique to emulated tracing.

\vheading{Inter-Device Dependencies} In distributed deep learning, collective communication operations are used to synchronize data across devices. To emulate them accurately, \sysname captures the full lifecycle of these collectives. Each worker initializes a communicator using an API like \texttt{ncclCommInitRank}, which assigns ranks and defines the communication topology for operations like \texttt{ncclAllReduce}. This setup enables \sysname to accurately track data dependencies and the role of each worker device within the collective operation.

Once initialized, these communicators facilitate data transfers that often run on dedicated streams to achieve compute-communication overlap. For example, in \texttt{ncclAllReduce}, each device concurrently performs compute tasks on one stream while executing collective communication on another. \sysname models this behavior by simultaneously tracking communication and compute streams, enabling it to accurately capture blocking dependencies and the resulting overlap between computation and data transfer. For more complex parallelism patterns, such as 3D parallelism, \sysname tracks multiple communicators operating across different workload dimensions --- assigning unique identifiers to each one and logging associated events in the trace. Just like compute operations, there is no need to actually share data between worker CPU processes since the control flow does not depend on the result of the collective; this obviates the need for IPC and synchronization in the emulator.

This approach works particularly well for DLT workloads due to their predictable, repetitive nature. The training loop typically executes the same sequence of operations repeatedly in each iteration, with control-flow decisions rarely depending on specific numerical results from device computation. By exploiting this characteristic while carefully maintaining device state, \sysname can accurately model workloads without executing device operations. The emulator produces detailed traces that capture the full complexity of device interactions while remaining lightweight and efficient.

\subsection{Trace Collection and Analysis}

\vheading{Worker Trace Generation} \sysname captures detailed execution traces for each worker in the distributed training job. Rather than just logging API calls, we maintain rich context about each operation. For compute kernels, we record essential metadata including input/output tensor shapes, data types, and memory layouts --- information critical for runtime prediction. For instance, in a transformer layer's attention computation, we track matrix dimensions and sparsity patterns that significantly impact performance.

Each trace entry also includes precise timing of CPU-side operations between kernel launches, capturing essential host overhead and dispatch latency. We achieve this by measuring wall-clock deltas between API calls during emulation. This is particularly important when operating with state-of-the-art devices like NVIDIA H100s, where dispatch overhead can be significant, especially in workloads with many small kernels.

\vheading{Trace Collation} A key challenge in analyzing distributed training workloads is reconstructing the global execution pattern from individual worker traces. While individual workers are aware of the number of participants in a collective operation, they do not have visibility into \textit{which} workers are involved or the topology of the communication graph. The collator identifies collective operations (like \texttt{ncclAllReduce}) and matches them across workers using communicator IDs and sequence numbers. This allows us to reconstruct and model the full communication pattern faithfully.

\label{sec:deduplication}
\vheading{Optimization: Worker Deduplication} An insight that enables \sysname to efficiently scale to large distributed workloads is that many workers in DLT perform identical work. For instance, in data parallel training, each worker executes the same computation on different data shards. We exploit this pattern through dynamic worker deduplication. During the first training iteration, we profile all workers to establish operation patterns. We compute rolling hashes of operation sequences, allowing us to identify workers performing redundant computation. Upon detecting duplicates, we terminate redundant workers and continue profiling only the unique ranks. The trace collator later reconstructs the full execution pattern using these profiled ranks.

This optimization is particularly effective for large-scale training jobs. For example, in a 64-GPU job with 8-way TP and 8-way DP, we only need to profile a single worker since tensor and data parallel workers exhibit identical behavior.


\subsection{End-to-end Simulator}

The emulator trace contains metadata for each operation but lacks execution times, since operations are emulated and not dispatched to actual hardware. To produce an end-to-end performance estimate from this trace, the simulator pipeline i) predicts per-operation runtimes from metadata recorded during emulation, and ii) conducts a discrete-event simulation of cluster behavior.

\vheading{Kernel Runtime Estimation} \sysname’s kernel runtime estimators are pluggable components that estimate latency and bandwidth for \textit{individual} compute and collective operations. Users can provide any runtime estimator of their choosing for any kernel type (eg. Habitat~\cite{geoffrey2021habitat}, GPU-Mangrove~\cite{braun2020mangrove}, static-analysis based approaches~\cite{alavani2018predstatic}).

By default, \sysname uses random forest regressors trained on profiling data from kernel microbenchmarks, similar to prior approaches \cite{zhu2020daydream, vidur}. For collective operations, the reference estimators leverage profiling data of intra-host and inter-host link characteristics, considering varying data sizes and the topology of participating devices. Please refer to Appendix \ref{appendix:predictions} for more details.

To facilitate easy onboarding of new operations, \sysname offers a transparent \textit{profiling} mode that dispatches operations on real hardware (rather than emulating them), logging each operation’s arguments and observed runtime. This enables us to progressively build and integrate prediction models from production workloads.


\vheading{Resource Model} The simulator models both host and accelerator resources. Each host machine is represented by a dispatch queue that processes operations and manages device interactions. Each accelerator manages multiple execution streams and models concurrent operation processing. Synchronization operations like \texttt{cudaDeviceSynchronize} and \texttt{cudaStreamWaitEvent} are modeled using blocking waits in the corresponding streams. Host-side computation and launch overheads are also modeled as blocking operations in the dispatch queue using measurements from the emulation phase. \sysname shares its overall discrete-event simulation approach with prior work \cite{zhu2020daydream, santhanam2021distir} - however, we can capture fine-grained dependencies at the CUDA API granularity owing to the detailed traces collected via emulation. For more details, please refer to Algorithm \ref{alg:simulator}, Appendix \ref{appendix:simulator_core}.

\vheading{Network Model} Network operations are implemented using a global waitmap where participating devices register themselves, blocking their respective streams until all workers join. This waitmap can capture pipeline bubbles and effectively model compute-comms overlap --- data dependencies manifest as stalls on the corresponding accelerator stream, while any concurrent compute streams can proceed to the next event (\autoref{fig:traceflow}). This behavior is described in more detail in Algorithm \ref{alg:waitmaps}, Appendix \ref{appendix:simulator_core}.

After all participants join, the on-the-wire duration of each collective operation is a black-box prediction from the corresponding kernel runtime estimator. This abstracts any topology-dependent runtime effects into a single discrete event that is separate from other dataflow dependencies in the simulator, allowing network operations to be modeled in isolation. This allows users to choose between profiled collective data from their target cluster (nccl-tests), or network simulators like ASTRA-sim \cite{won2023astra}.

By reproducing the behavior of accelerator primitives, the simulator provides an accurate representation of cluster behavior. Operation-level modeling ensures we can capture fine-grained behavior while remaining general; new computational optimizations can be captured without additional effort. While our implementation targets CUDA devices, the design generalizes to other accelerators.
